\documentclass[11pt]{article}
\usepackage[a4paper,left=16mm,right=16mm,top=13mm,bottom=13mm]{geometry}
\usepackage[T1]{fontenc}
\usepackage{lmodern,microtype,amsmath,amssymb,mathtools}
\usepackage[hidelinks]{hyperref}
\hypersetup{pdftitle={Explicit coordinate functions for a dominant map to the Klein cubic},pdfauthor={Research notes}}
\pagestyle{empty}
\setlength{\parindent}{0pt}
\setlength{\parskip}{2.5pt}
\setlength{\abovedisplayskip}{5pt}
\setlength{\belowdisplayskip}{5pt}
\setlength{\abovedisplayshortskip}{3pt}
\setlength{\belowdisplayshortskip}{3pt}
\DeclareMathOperator{\adj}{adj}
\DeclareMathOperator{\tr}{tr}
\DeclareMathOperator{\Pf}{Pf}
\newcommand{\PP}{\mathbf P}
\begin{document}
\begin{center}
{\Large\bfseries Explicit coordinates for the Klein cubic map}\\[3pt]
$\displaystyle \Phi:\PP^{19}\dashrightarrow X=\bigl\{\textstyle\sum_{j\bmod5}y_j^2y_{j+1}=0\bigr\}\subset\PP^4.$
\end{center}
\vspace{-4pt}
Use coordinates $a=(a_{ij})_{0\le i\le3,\,0\le j\le4}$ on four copies of the honest Klein representation
$A=\mathbf C^5$, with diagonal $G=\mathrm{PSL}_2(\mathbf F_{11})$-action. All indices start at zero.
The following formulas specify five homogeneous polynomials $R_j(a)$; the map is $[R_0:\cdots:R_4]$.

\textbf{1. Fixed matrices and bases.}
In the coordinate basis $e_0,\ldots,e_5$ of the Pfaffian six-space, set $B_i=M(a_{i0},\ldots,a_{i4})$, where
\[
M(y)=\begin{pmatrix}
0&-y_1&-y_3&-y_2&-y_0&-y_4\\
y_1&0&-y_0&0&0&y_3\\
y_3&y_0&0&0&-y_2&0\\
y_2&0&0&0&y_4&-y_1\\
y_0&0&y_2&-y_4&0&0\\
y_4&-y_3&0&y_1&0&0
\end{pmatrix},\qquad
L_i{}_{h,rst}=\delta_{hr}(B_i)_{st}-\delta_{hs}(B_i)_{rt}+\delta_{ht}(B_i)_{rs}.
\]
The $6\times20$ matrix $L_i$ uses row basis $e_0,\ldots,e_5$; column $rst$ means $e_r\wedge e_s\wedge e_t$.
Order columns by
\[
\mathcal J=(012,013,014,015,023,024,025,034,035,045,
123,124,125,134,135,145,234,235,245,345).
\]
Stack $L_0,L_1,L_2$ in that order and write $\mathcal C=(N\mid r\mid s)$, where $N$ has the first
18 columns. Put $\Delta=\det N$. For those columns $J$, define
\[
u_J=-\det N[J\leftarrow r],\qquad v_J=-\det N[J\leftarrow s],\qquad
(u_{245},u_{345})=(\Delta,0),\quad(v_{245},v_{345})=(0,\Delta).
\]
The notation $N[J\leftarrow r]$ means replace column $J$ by $r$, without changing its position.

\textbf{2. Three scalar coefficients.}
For a twenty-vector $w$ indexed by $\mathcal J$, put
\[
\pi(w)=w_{012}w_{034}-w_{013}w_{024}+w_{014}w_{023},\qquad
q_0=\pi(u),\quad q_1=\pi(u+v)-\pi(u)-\pi(v),\quad q_2=\pi(v).
\]
No roots of $q_2z^2+q_1z+q_0$ are needed. All subsequent calculations are ordinary scalar arithmetic.

\textbf{3. Four pairs of $3\times3$ matrices.}
For $\epsilon=0,1$, let $w^{(0)}=u$, $w^{(1)}=v$, and form the $3\times4$ matrices
\[
D^{(\epsilon)}_{hi}=\sum_{k=0}^5\sum_{J\in\mathcal J}(B_i)_{hk}(L_3)_{kJ}w^{(\epsilon)}_J
\quad(0\le h\le2,\ 0\le i\le3).
\]
Let $E_i$ and $H_i$ be $D^{(0)}$ and $D^{(1)}$ with column $i$ deleted, retaining the other columns
in increasing order. Define
\[
(c_{i0},c_{i1},c_{i2},c_{i3})=(-1)^i
\bigl(\det E_i,\ \tr(\adj(E_i)H_i),\ \tr(\adj(H_i)E_i),\ \det H_i\bigr).
\]
Here $\adj$ is the transposed cofactor matrix; thus every $c_{im}$ is a specified polynomial.

\textbf{4. The five coordinate functions.}
For $j=0,\ldots,4$, set
\[
\begin{aligned}
P_j&=\sum_{i=0}^3 a_{ij}\bigl(q_2^2c_{i0}-q_0q_2c_{i2}+q_0q_1c_{i3}\bigr),\\
Q_j&=\sum_{i=0}^3 a_{ij}\bigl(q_2^2c_{i1}-q_1q_2c_{i2}+(q_1^2-q_0q_2)c_{i3}\bigr),\\
S&=\sum_{m=0}^4 Q_m^2Q_{m+1},\qquad
T=\sum_{m=0}^4(2Q_mQ_{m+1}+Q_{m-1}^2)P_m\quad\text{(subscripts modulo 5).}
\end{aligned}
\]
\[
\boxed{\quad R_j=q_2SP_j+(q_1S-q_2T)Q_j\quad(j=0,1,2,3,4),\qquad
\Phi(a)=[R_0:R_1:R_2:R_3:R_4].\quad}
\]
On $R_0\ne0$ the four affine rational functions are $R_1/R_0,\ldots,R_4/R_0$.
The $R_j$ have common degree $568$ before cancelling common factors.

\textbf{Check.}
For the input with rows $(1,1,0,0,1)$, $(0,1,-1,1,0)$, $(-1,-1,1,-1,-1)$,
$(-1,1,-1,0,0)$, one obtains $\Delta=q_2=1$, $(q_0,q_1)=(1,3)$ and
$R=(-3,2,-2,1,-2)$. At this input,
\[
\det\frac{\partial(R_1/R_0,\,R_2/R_0,\,R_3/R_0)}{\partial(a_{00},\,a_{02},\,a_{10})}
=-\frac{221}{27}\ne0.
\]
The residual-secant identity gives $\sum_jR_j^2R_{j+1}=0$; the intrinsic construction gives
$G$-equivariance. The nonzero minor proves dominance.

\vfill
{\footnotesize Pfaffian model: Tschinkel--Zhang, arXiv:2409.08392, \S4.
The earlier four-copy map, with all operations written in coordinates.\par}
\end{document}
